\chapter{Introducción}

\section{Objetivos de la práctica}
El objetivo de esta práctica es la configuración 
de un servidor web seguro (Apache con HTTPS) 
mediante el uso de certificados digitales 
firmados por la Autoridad de Certificación 
(CA) de la asignatura. Para alcanzar este 
propósito se plantean los siguientes objetivos:
\begin{itemize}
    \item Generar y firmar un certificado digital para servidor web que valide identidades múltiples (\texttt{ejemplo.com} y \texttt{www.ejemplo.com}) utilizando la extensión \texttt{subjectAltName} de X.509 v3.
    \item Manipular la resolución de nombres de dominio de forma local modificando el archivo \texttt{/etc/hosts}.
    \item Incorporar certificados raíz personalizados a la lista de autoridades de confianza del sistema Linux.
    \item Instalar, configurar y administrar de forma segura un servidor web Apache.
    \item Comprobar y verificar la robustez de las conexiones seguras TLS a través de herramientas de diagnóstico (\texttt{openssl s\_client}) y navegadores convencionales (Mozilla Firefox).
\end{itemize}

\section{Fundamentos teóricos}
Para la realización de esta práctica, se aplican diversos conceptos teóricos clave de la ciberseguridad y la gestión de redes:
\begin{itemize}
    \item \textbf{Certificados Digitales y X.509:} Un certificado digital vincula criptográficamente una clave pública a la identidad de un servidor, dominio o usuario, y es validado por la firma de una Autoridad de Certificación. El estándar \textbf{X.509 v3} permite incluir extensiones avanzadas; una de las más relevantes para servidores web modernos es el \textit{Subject Alternative Name} (SAN), imprescindible para que un único servicio (y certificado) responda e identifique de forma legítima a múltiples nombres DNS.
    \item \textbf{HTTPS y TLS (Transport Layer Security):} HTTPS no es más que HTTP sobre una capa de seguridad TLS. TLS asegura las comunicaciones aportando confidencialidad (cifrado de datos transmitidos), integridad (protección contra alteraciones de datos en tránsito) y autenticación del lado del servidor (para evitar ataques de intermediarios).
    \item \textbf{Resolución Local de Dominios:} Antes de interrogar a servidores DNS públicos, los sistemas operativos consultan el archivo \texttt{/etc/hosts} local para la resolución de nombres. Operar sobre este archivo permite emular el control de un dominio legítimo dentro de un entorno de laboratorio o máquina virtual aislada.
\end{itemize}

\section{Enunciado de la práctica}
El desarrollo del ejercicio y de la presente memoria seguirán el guion de actividades propuesto en el enunciado:

\begin{enumerate}
    \item Creación del certificado del servidor \texttt{www.ejemplo.com} (y \texttt{ejemplo.com}) firmado con la autoridad de la asignatura. Asegurarse de que el certificado incorpora correctamente los nombres \texttt{www.ejemplo.com} y \texttt{ejemplo.com} en la extensión \texttt{subjectAltName}.
    \item Modificación del fichero \texttt{/etc/hosts} de la máquina virtual y verificación de los alias \texttt{www.ejemplo.com} y \texttt{ejemplo.com}.
    \item Incorporación del certificado RAÍZ de la asignatura a la lista de autoridades reconocidas en el servidor Linux de la máquina virtual.
    \item Configuración del servidor Apache. Descripción detallada de los ficheros modificados, certificados de cliente y CA.
    \item Instalación y arranque del servidor Apache. Instalación de paquetes, ejecución de comandos de administración del servicio, etc. Revisión continua del fichero \texttt{error\_log} para detectar warnings y errores; y del \texttt{access\_log} para comprobar accesos.
    \item Creación y configuración del fichero \texttt{index.html}. Permisos y configuración adicional necesaria.
    \item Conexión desde \texttt{openssl s\_client} a \texttt{www.ejemplo.com:443} y comparación de las primeras líneas de la conexión con la conexión a \texttt{www.ulpgc.es:443}; asegurarse de que no hay problema en la identificación del servidor y los certificados que presenta.
    \item Conexión desde el navegador local. Firefox debe acceder sin mostrar ningún error. No permitir la conexión si aparecen errores del tipo ``entiendo los problemas...''.
    \item Capturas de pantalla del Firefox (cuatro como mínimo) para ambos dominios, una mostrando el contenido del fichero \texttt{index.html} y otra con las propiedades de la conexión (Herramientas / Información de la página / Seguridad en Firefox o bien pulsar en el ``candado'' de la conexión segura en Firefox, más información...), mostrando la versión de TLS y el tipo de cifrado obtenido en cada una de las conexiones.
\end{enumerate}

\section{Nota sobre el \emph{prompt} personalizado}
En ocasiones se ha omitido la cabecera de los codigos insertados en la memoria,
esto no es un error, ni que se hayan obtenido de manera externa, sino que al mostrar el contenido
de los certificados, claves o configuraciones, se ha optado por mostrar únicamente el contenido relevante, 
sin incluir el prompt personalizado que se muestra en la terminal de la máquina virtual. 